\section{Appendice computazionale: costruzione e audit numerico}
\label{app:computational}

Questa appendice fornisce una famiglia esplicita di grafi raffinati e un
protocollo computazionale. I calcoli sono da intendersi come
verifiche numeriche di proprietà teoricamente formulate; non sostituiscono
le dimostrazioni analitiche.

\subsection{Famiglia di griglie e embedding}

Sia $M$ una varietà riemanniana spaziale compatta di dimensione $d$,
con atlante finito e metrica $h\in C^5$. Per semplicità espositiva,
si consideri inizialmente il toro piatto
$M=\mathbb T^d=\mathbb R^d/\mathbb Z^d$. Per
$a=1/n$, $n\in\mathbb N$, definiamo
\[
V_a=(\mathbb Z/n\mathbb Z)^d,
\qquad
X_a(k)=k/n\pmod{\mathbb Z^d}.
\]
Connettiamo $k$ a $k\pm e_i$ e poniamo peso unitario su ciascun link.
La distanza dei link è $a$ e la valenza è $2d$ indipendentemente da
$a$. Il Laplaciano discreto è
\[
\Delta_a f(k)=\frac1{a^2}\sum_{i=1}^d
\bigl[f(k+e_i)+f(k-e_i)-2f(k)\bigr].
\]
Per una varietà generale si costruiscono triangolazioni quasi-uniformi
con diametro massimo $a$, si scelgono coordinate normali locali e si
impongono pesi consistenti con il tensore metrico. Le ipotesi di
shape-regularity e quasi-uniformità devono essere uniformi nella
famiglia.

\subsection{Audit dei momenti}

Per ogni vertice e ogni direzione si calcolano
\[
M_{1,a}^\nu=\frac1{a^2}\sum_yw_a\xi_{xy}^\nu,
\qquad
M_{2,a}^{\mu\nu}=\frac1{2a^2}\sum_yw_a\xi_{xy}^\mu\xi_{xy}^\nu.
\]
L'audit deve registrare
\[
E_{1,a}=\max_x\|M_{1,a}(x)-b(x)\|,
\qquad
E_{2,a}=\max_x\|M_{2,a}(x)-h^{-1}(x)\|,
\]
insieme al massimo del terzo momento e del quarto momento normalizzato.
Una regressione di $\log E$ contro $\log a$ stima soltanto l'ordine
osservato sul campione di scale testato.

\subsection{Ricostruzione metrica e stabilità inversa}

Sia $\mathcal R_a$ l'operatore di ricostruzione che associa ai momenti
la matrice metrica stimata $\widehat h_a$. Per ottenere convergenza
metrica occorre assumere, oltre alla convergenza dei momenti,
\[
\|\mathcal R_a(M_{2,a})-\mathcal R_a(h^{-1})\|
\leq L\|M_{2,a}-h^{-1}\|,
\]
con $L$ indipendente da $a$. Se inoltre
\[
\|M_{2,a}-h^{-1}\|\leq Ca^2,
\]
allora
\[
\|\widehat h_a-h\|\leq LCa^2
\]
quando l'inversione matriciale resta uniformemente lontana dalla
singolarità. In particolare, se gli autovalori di $h$ sono compresi in
$[\lambda_-,\lambda_+]$ con $\lambda_->0$ e l'errore è inferiore a
$\lambda_-/2$, l'inversione è uniformemente stabile.

\subsection{Test di stabilità}

Per un problema ellittico discreto $A_au_a=f_a$, si deve controllare
una stima uniforme
\[
\|u_a\|_{X_a}\leq C\|f_a\|_{Y_a},
\]
con $C$ indipendente da $a$. Per il settore lorentziano, tale test
non è sufficiente: occorre una famiglia di operatori evolutivi
$U_a(t,s)$ e una stima di energia
\[
\|u_a(t)\|_{E_a}\leq C_T\left(\|u_a(0)\|_{E_a}
+\int_0^t\|F_a(\tau)\|\,\dd\tau\right),
\qquad 0\leq t\leq T,
\]
con $C_T$ indipendente da $a$.

\subsection{Audit Schwarzschild}

Nella regione esterna si confrontano le componenti ricostruite con
\[
g_{tt}^{\rm Schw}=-(1-2M/r),
\qquad
g_{rr}^{\rm Schw}=(1-2M/r)^{-1}.
\]
La condizione geometrica forte non è soltanto l'accordo di due
componenti: occorre controllare il tensore di Einstein discreto,
la simmetria sferica, la convergenza dei vincoli e la stabilità rispetto
alla regione di fitting. Un risultato numerico è compatibile con
Schwarzschild soltanto se l'errore tende a zero sotto raffinamento e
se le ipotesi di vuoto e di identificazione del settore Einstein sono
soddisfatte.

\subsection{Pseudo-codice riproducibile}

\begin{verbatim}
for n in refinement_levels:
    a = 1.0/n
    G = build_quasi_uniform_graph(M, a)
    X = build_embedding(G, M)
    moments = compute_moments(G, X)
    metric = reconstruct_metric(moments)
    residuals = compute_moment_residuals(moments, metric)
    stability = solve_discrete_problem(G, boundary_data)
    store(a, residuals, stability, metric)
fit_loglog(error, a)
compare_with_schwarzschild(metric, M)
\end{verbatim}

\subsection{Limiti dell'audit}

Senza coordinate effettive, pesi, precisione numerica, condizioni al
bordo, dati grezzi e definizione completa dell'operatore lorentziano,
non è possibile certificare una convergenza numerica completa. I test
forniti nel corpo principale restano quindi evidenza computazionale
riportata, mentre i teoremi richiedono le ipotesi analitiche esplicitate.
